\documentclass[11pt]{article}
\usepackage{coursenotes}
\begin{document}
\notesheader{10}{Panel data and difference-in-differences}{Untreated units over time as the comparison}

\begin{decision}
A change was rolled out to some stores, clubs, or cities and not others, on one date, and nobody randomized which. The firm wants
to know what it did before extending it. The comparison now comes from untreated units' changes over time. That requires an
assumption about what treated units would have done without the change, which no data can verify, so the analysis must also say
how far that assumption could fail before the decision changes.
\end{decision}

\begin{goals}
  \item explain why before--after comparisons fail, including regression to the mean;
  \item use fixed effects and the within transformation, and say what they cannot remove;
  \item identify the ATT by difference-in-differences under parallel trends, with a proof;
  \item show that TWFE with one adoption date equals the DiD, and read an event study;
  \item compute a breakdown drift for trend sensitivity, and cluster standard errors correctly;
  \item recognize why staggered adoption breaks the simple TWFE regression.
\end{goals}

%=====================================================================================
\sectionone

\subsection{Setup and notation}

Units $i$ are observed in periods $t$. Index potential outcomes by \emph{adoption date}: $Y_{it}(g)$ is unit $i$'s outcome in period $t$
if it first adopts in period $g$, and $g = \infty$ means never. With one adoption date $t_0$, treated units ($T_i = 1$) have
$Y_{it}(t_0)$ observed and comparison units ($T_i = 0$) have $Y_{it}(\infty)$. Write $Y_{it}(1) \equiv Y_{it}(t_0)$ and
$Y_{it}(0) \equiv Y_{it}(\infty)$ for short. The target is the effect on the treated after adoption,
$\ATT = \E[Y_{\text{post}}(1) - Y_{\text{post}}(0) \mid T = 1]$, whose difficulty is $\E[Y_{\text{post}}(0) \mid T = 1]$.

Nothing in the design made adoption random. The assumptions below replace the coin of Meetings 1--6 and the recorded rule of
Meetings 7--8. They are claims about untreated outcomes, supported by facts about \emph{why} units adopted when they did, and
they can be checked only in part. Identification here rests on assumptions, and the analysis should say so.

\subsection{Why before--after fails}

A treated unit's before--after change attributes to the treatment everything else that changed: trends, seasons, other events,
and regression to the mean.

\begin{prop}{Before--after with selection on a bad period}{rtm}
Let $Y_{it}(0) = a_i + \lambda_t + u_{it}$ with transitory shocks $u_{it}$ independent over time with mean zero. If units are treated
because their pre-period outcome was low, the before--after change among treated units is
\[
  \E[Y_{\text{post}} - Y_{\text{pre}} \mid T = 1] = \ATT + (\lambda_{\text{post}} - \lambda_{\text{pre}}) - \E[u_{\text{pre}} \mid T = 1],
\]
and the last term is positive.
\end{prop}
\begin{proof}
Substitute the model: the change is $\ATT + \lambda_{\text{post}} - \lambda_{\text{pre}} + \E[u_{\text{post}} - u_{\text{pre}} \mid T=1]$. Selection uses
$Y_{\text{pre}}$ only, so $\E[u_{\text{post}} \mid T=1] = 0$, while selecting low $Y_{\text{pre}}$ makes $\E[u_{\text{pre}} \mid T=1] < 0$.
\end{proof}

``Fix the worst-performing stores'' builds regression to the mean into the rollout. A comparison group fixes the trend term but not this
one unless it was selected the same way.

\subsection{Fixed effects}

The two-way fixed-effects (TWFE) model $Y_{it} = \alpha_i + \lambda_t + \beta D_{it} + \varepsilon_{it}$ gives each unit a level $\alpha_i$ (absorbing every
stable characteristic, measured or not) and each period a shock $\lambda_t$ common to all units. Confounders that change over time
\emph{differently} across units (a neighborhood gentrifying, a manager who adopts and works harder) are not removed.

\begin{prop}{TWFE by the within transformation}{within}
In a balanced panel, with $\ddot W_{it} = W_{it} - \bar W_{i\cdot} - \bar W_{\cdot t} + \bar W_{\cdot\cdot}$, the least-squares coefficient on $D$ is
$\hat\beta = \sum_{i,t}\ddot D_{it}Y_{it}/\sum_{i,t}\ddot D_{it}^2$.
\end{prop}
\begin{proof}
Unit and period indicators span the additive functions $a_i + b_t$; in a balanced panel the residual from projecting on them is the
two-way demeaned variable. The Frisch--Waugh--Lovell theorem from Meeting~8 finishes it: the coefficient on $D$ is the slope of $Y$
on the part of $D$ left after projecting out the other regressors. In an unbalanced panel the one-pass demeaning is not the
projection, and the regression with both sets of indicators must be run directly.
\end{proof}

Only variation within units over time identifies $\beta$; variables that never change within a unit drop out.

\subsection{Difference-in-differences}

\begin{assum}{Parallel trends}{pt}
$\E[Y_{\text{post}}(0) - Y_{\text{pre}}(0) \mid T=1] = \E[Y_{\text{post}}(0) - Y_{\text{pre}}(0) \mid T=0]$.
\end{assum}
\begin{assum}{No anticipation and no spillovers}{noant}
$Y_{\text{pre}}(1) = Y_{\text{pre}}(0)$ for treated units, and comparison units' outcomes do not depend on whether treated units are treated.
\end{assum}

Parallel trends restricts \emph{changes}, not levels: the groups may differ by any amount provided the difference would have stayed constant.
It is untestable, because the treated group's post-period $Y(0)$ is never seen. No anticipation fixes what ``pre'' means. If behavior changed
when adoption was announced, the pre-period must end before that.

\begin{prop}{DiD identifies the ATT}{did}
Under Assumptions~\ref{asm:pt}--\ref{asm:noant},
\[
  \ATT = \{\E[Y_{\text{post}} \mid T=1] - \E[Y_{\text{pre}} \mid T=1]\} - \{\E[Y_{\text{post}} \mid T=0] - \E[Y_{\text{pre}} \mid T=0]\}.
\]
\end{prop}
\begin{proof}
For treated units the observed means are $\E[Y_{\text{post}}(1) \mid T=1]$ and, by no anticipation, $\E[Y_{\text{pre}}(0) \mid T=1]$; for comparison units
both are of $Y(0)$, free of spillovers. Replace the comparison group's change by the treated group's counterfactual change (parallel trends);
the $Y_{\text{pre}}(0)$ terms cancel.
\end{proof}

The before--after change keeps the common \emph{trend}; the post-period cross-section keeps the \emph{level gap}, which is the selection bias of
Meeting~1, $\E[Y(0) \mid T=1] - \E[Y(0) \mid T=0]$, measured before treatment. The DiD removes both.

\begin{prop}{Parallel trends depends on the scale}{scale}
If parallel trends holds in levels with a common counterfactual change $\Delta \ne 0$ and the groups' pre-period means differ, it fails in
logs (approximately, in percentage changes).
\end{prop}
\begin{proof}
The implied counterfactual percentage changes are $\Delta/\mu_1$ and $\Delta/\mu_0$, unequal when $\mu_1 \ne \mu_0$.
\end{proof}
Choose the scale on which common shocks plausibly act (multiplicative shocks suggest logs), and report the other if the decision is close.

\subsection{TWFE, event studies, and pre-trends}

\begin{prop}{With one adoption date, TWFE is the DiD}{twfedid}
In a balanced panel with $n_1$ treated and $n_0$ comparison units, $\mathcal T_0$ pre- and $\mathcal T_1$ post-periods, and
$\Delta_i = \bar Y_{i,\text{post}} - \bar Y_{i,\text{pre}}$, the TWFE coefficient is $\hat\beta = \bar\Delta_{T=1} - \bar\Delta_{T=0}$.
\end{prop}
\begin{proof}
Here $\ddot D_{it} = (T_i - \bar T)(P_t - \bar P)$ with $P_t$ the post indicator. Then $\sum_t(P_t - \bar P)Y_{it} = (\mathcal T_0\mathcal T_1/\mathcal T)\Delta_i$ and
$\sum_i(T_i - \bar T)\Delta_i = (n_1n_0/N)(\bar\Delta_{T=1} - \bar\Delta_{T=0})$; the denominator factors the same way.
\end{proof}

The regression adds convenience (many periods, standard errors), not identification. An \emph{event study} replaces $D_{it}$ by indicators
$T_i\ind{t - t_0 = k}$ for each period relative to adoption except a base $k = -1$. In a balanced panel each coefficient is
$\hat\beta_k = (\bar Y_{1k} - \bar Y_{0k}) - (\bar Y_{1,-1} - \bar Y_{0,-1})$: post-period coefficients trace the effect over time; pre-period
coefficients are placebo checks, near zero under parallel trends and no anticipation. The base period's coefficient is zero by definition and
has no error bar.

Flat pre-period coefficients are weak evidence: with few, noisy pre-periods a test can miss a drift large enough to explain the whole
``effect''; reporting only analyses that pass the test biases what gets reported; and a shock that coincides with adoption (a local campaign)
leaves pre-trends untouched.

\begin{prop}{Linear differential trend and the breakdown drift}{drift}
If the counterfactual gap drifts linearly, $\E[Y_t(0) \mid T=1] - \E[Y_t(0) \mid T=0] = a + \delta t$, a DiD comparing post periods with mean
$\bar t_{\text{post}}$ to pre periods with mean $\bar t_{\text{pre}}$ converges to $\ATT + \delta(\bar t_{\text{post}} - \bar t_{\text{pre}})$. Against a break-even $c$,
the decision changes at the \emph{breakdown drift} $\delta^\ast = (\text{DiD} - c)/(\bar t_{\text{post}} - \bar t_{\text{pre}})$.
\end{prop}
\begin{proof}
The DiD is the post-period average gap minus the pre-period average gap; the counterfactual part of the difference is
$\delta(\bar t_{\text{post}} - \bar t_{\text{pre}})$.
\end{proof}

Report $\delta^\ast$ in the decision's units, compare it with the drift the pre-period coefficients can rule out, and argue from the institution
whether a drift that large is plausible. Later event-time coefficients and distant bases are more exposed to drift.

\subsection{Inference}

A unit's outcomes are correlated over time and treatment is constant within long blocks, so unit-period standard errors are too small, often
by a large factor.

\begin{prop}{Collapsing to one change per unit}{collapse}
If units are independent, $\bar\Delta_{T=1} - \bar\Delta_{T=0}$ has $\SE = \sqrt{s_1^2/n_1 + s_0^2/n_0}$ with $s_d$ the standard deviation of $\Delta_i$ in
group $d$, whatever the correlation of a unit's outcomes over time.
\end{prop}
\begin{proof}
Each unit contributes one $\Delta_i$; within-unit correlation affects $\Var(\Delta_i)$, not the independence of $\Delta_i$ across units.
\end{proof}

With one adoption date this is asymptotically the TWFE SE clustered by unit. Cluster at the level at which treatment was assigned and shocks are
shared (by region if regions adopted). With few clusters use $t$ critical values with cluster-based degrees of freedom or the wild cluster
bootstrap.

\subsection{A warning about staggered adoption}

When units adopt at different dates, TWFE with a single $D_{it}$ still runs but averages many two-by-two comparisons, some using
\emph{already-treated} units as controls for later adopters. If effects grow over time, those comparisons subtract part of the effect, some get
negative weight, and the estimate can be badly biased, even of the wrong sign. Use estimators that compare only with not-yet-treated or
never-treated units, and report effects by cohort.

\begin{pitfall}[DiD errors]
Before--after on units chosen for a bad period; controls changed by the treatment; unit-period SEs; a base period inside an anticipation window;
reading flat pre-trends as proof; TWFE with staggered adoption.
\end{pitfall}

%=====================================================================================
\sectiontwo

\paragraph{Setting.} Yangsi Fitness introduced app-based class booking in 20 of its 50 clubs at the start of July; the other 30 kept phone and desk
booking. The outcome is average monthly visits per member. The system costs ¥30{,}000 per club a year; management will extend it if it raises
visits by at least $0.3$ per member per month.

\paragraph{Step 1: before--after.} App clubs: $8.0$ visits (January--June) and $8.9$ (July--December), a change of $0.9$, which includes the autumn
rise every club sees.

\paragraph{Step 2: the two-by-two DiD.}
\begin{center}
\begin{tabular}{lccc}
\toprule
 & Jan--Jun & Jul--Dec & Change \\
\midrule
App clubs (20) & 8.0 & 8.9 & $+0.9$ \\
Other clubs (30) & 7.5 & 7.9 & $+0.4$ \\
\midrule
Difference & 0.5 & 1.0 & DiD $= +0.5$ \\
\bottomrule
\end{tabular}
\end{center}
The other clubs' rise of $0.4$ stands in for the seasonal rise app clubs would have had. The effect is $0.5$ visits per member per month (about 6\% of
the counterfactual $8.4$), above $0.3$. App clubs were busier to begin with; parallel trends does not need equal levels.

\paragraph{Step 3: levels or logs.} In logs, $\ln(8.9/8.0) - \ln(7.9/7.5) = 0.055$; parallel trends in logs implies a counterfactual of
$8.0 \times 7.9/7.5 = 8.43$ and an effect of $0.47$ visits. The two scales give $0.50$ and $0.47$, and the decision is the same.

\paragraph{Step 4: event study.} Monthly TWFE with June as the base, SEs clustered by club (50 clusters):
\begin{center}
\begin{tikzpicture}
\begin{axis}[width=0.66\linewidth, height=5cm, xlabel={months relative to adoption}, ylabel={$\hat\beta_k$}, axis lines=left,
  xmin=-3.6, xmax=2.6, ymin=-0.25, ymax=0.8, xtick={-3,-2,-1,0,1,2}]
\addplot[only marks, mark=*, sbgreen, error bars/.cd, y dir=both, y explicit]
  coordinates {(-3,0.02) +- (0,0.118) (-2,-0.03) +- (0,0.118) (0,0.35) +- (0,0.118) (1,0.52) +- (0,0.118) (2,0.60) +- (0,0.118)};
\addplot[only marks, mark=o, extgrey] coordinates {(-1,0)};
\draw[dashed, extgrey] (axis cs:-3.6,0) -- (axis cs:2.6,0);
\draw[thick, dashed, warnred] (axis cs:-3.6,0.3) -- (axis cs:2.6,0.3);
\end{axis}
\end{tikzpicture}

\small Figure 1. Event-study coefficients for months $-3$ to $+2$ with 95\% intervals (SE $0.06$); base month $-1$ has none. Red: the $0.3$ break-even.
\end{center}
Pre-period coefficients ($0.02$ and $-0.03$) are small; the effect builds over three months as members learn the app.

\paragraph{Step 5: trend sensitivity.} Under a linear drift $\delta$, the pre-period coefficients would be $\beta_{-2} = -\delta$ and $\beta_{-3} = -2\delta$.
Fitting $\delta$ to them by least squares gives $\hat\delta = -(1 \times (-0.03) + 2 \times 0.02)/5 = -0.002$ with SE $0.06/\sqrt5 = 0.027$, so the 95\% interval for the drift is
$[-0.055, 0.051]$: a drift of about $0.05$ a month in either direction cannot be ruled out. The breakdown drift for the month-$+2$ coefficient (three months from the base) is
$\delta^\ast = (0.60 - 0.30)/3 = 0.10$, nearly twice what the pre-period allows. Even a drift of $0.055$ leaves $0.60 - 3 \times 0.055 = 0.435 > 0.3$.

\paragraph{Step 6: how clubs were chosen.} The 20 clubs had the most complaints about phone booking; their visits had been stable, which argues
against regression to the mean. Had they been chosen after a slump, the event study would show a dip before adoption.

\paragraph{Recommendation.} Extend the app. The remaining 30 clubs will adopt at different times, so evaluate the extension with an estimator
that compares each wave only with clubs not yet on the app, ideally with adoption dates randomized across clubs.

\begin{exercise}[computation]
Collapsing to one change per club, the 20 app clubs' changes $\Delta_i$ have SD $0.20$ and the 30 other clubs' SD $0.18$. Compute the SE of the DiD and a
95\% interval. Clubs belong to six regions with shared marketing; what would you do about that?
\end{exercise}
\answer{$\SE = \sqrt{0.20^2/20 + 0.18^2/30} = \sqrt{0.0020 + 0.0011} = 0.055$; interval $0.5 \pm 1.96 \times 0.055 = [0.39, 0.61]$, clear of the
break-even of $0.3$. But if regions share shocks, the clubs are not independent: cluster by region. With only 6 clusters, use a $t$ distribution with about
5 degrees of freedom or the wild cluster bootstrap, and expect a much wider interval that may reach the break-even.}

\begin{exercise}[judgment]
A regional manager wants to report the app's effect as the before--after change in the 20 clubs, ``because the other clubs are in different cities and not
comparable''. Respond.
\end{exercise}
\answer{Before--after includes the seasonal rise ($0.4$ here) and anything else that changed. DiD does not need comparable levels, only comparable changes
(parallel trends), which the event study's flat pre-period supports. If cities differ in seasonality, restrict the comparison to clubs in similar markets
or use conditional parallel trends, rather than dropping the comparison group.}

%=====================================================================================
\sectionthree

\begin{extension}[Conditional parallel trends, few clusters, changes-in-changes]
If parallel trends is plausible only within covariate groups, outcome-regression, IPW, and doubly robust DiD estimators compare changes within groups
(Sant'Anna and Zhao 2020). With few clusters, the wild cluster bootstrap gives better-sized tests. Changes-in-changes (Athey and Imbens 2006) replaces
parallel means with a stable mapping between quantiles.
\end{extension}

\begin{extension}[Triple differences and randomized adoption]
A triple difference adds a within-unit comparison group that should be unaffected (members of a class type not bookable by app). Randomizing adoption
dates turns DiD into a design-based comparison. When dates are assigned at random, comparing an early cohort's mean outcome with a later
cohort's, before the later one adopts, is unbiased for the effect under no anticipation, with no parallel-trends assumption.
\end{extension}

\subsection*{Annotated reading}
\begin{readinglist}
\reading{Angrist, Joshua D., and J\"orn-Steffen Pischke. 2009. \emph{Mostly Harmless Econometrics: An Empiricist's
Companion}. Princeton: Princeton University Press.}{DiD and fixed effects at textbook pace}{Ch.\ 5}{2}
\reading{Bertrand, Marianne, Esther Duflo, and Sendhil Mullainathan. 2004. ``How Much Should We Trust
Differences-in-Differences Estimates?'' \emph{Quarterly Journal of Economics} 119 (1): 249--275.
\url{https://doi.org/10.1162/003355304772839588}}{Why unit-period SEs mislead}{All}{1}
\reading{Roth, Jonathan. 2022. ``Pretest with Caution: Event-Study Estimates after Testing for Parallel Trends.''
\emph{American Economic Review: Insights} 4 (3): 305--322. \url{https://doi.org/10.1257/aeri.20210236}}{What a passed
pre-trend test does to the estimate}{All}{2}
\reading{Rambachan, Ashesh, and Jonathan Roth. 2023. ``A More Credible Approach to Parallel Trends.'' \emph{Review of
Economic Studies} 90 (5): 2555--2591. \url{https://doi.org/10.1093/restud/rdad018}}{Sensitivity beyond linear
drift}{Sections 1--3}{3}
\reading{Roth, Jonathan, Pedro H. C. Sant'Anna, Alyssa Bilinski, and John Poe. 2023. ``What's Trending in
Difference-in-Differences? A Synthesis of the Recent Econometrics Literature.'' \emph{Journal of Econometrics} 235 (2):
2218--2244. \url{https://doi.org/10.1016/j.jeconom.2023.03.008}}{A survey of modern DiD}{All}{2}
\reading{Sant'Anna, Pedro H. C., and Jun Zhao. 2020. ``Doubly Robust Difference-in-Differences Estimators.''
\emph{Journal of Econometrics} 219 (1): 101--122. \url{https://doi.org/10.1016/j.jeconom.2020.06.003}}{Outcome-regression,
IPW, and doubly robust DiD under conditional parallel trends}{Introduction}{3}
\reading{Athey, Susan, and Guido W. Imbens. 2006. ``Identification and Inference in Nonlinear Difference-in-Differences
Models.'' \emph{Econometrica} 74 (2): 431--497. \url{https://doi.org/10.1111/j.1468-0262.2006.00668.x}}{Changes-in-changes:
a stable mapping between quantiles in place of parallel means}{Introduction}{3}
\reading{Handout H6, \emph{Staggered adoption}}{Cohorts, clean comparisons, and the TWFE decomposition}{All}{2}
\end{readinglist}

\end{document}
